\chapter{Results and Discussion}
\label{chap:results}

% NOTE: \S\ref{sec:results:10bar} is pre-filled with real numbers from
% \texttt{results/scoreboard.csv}. All remaining sections are
% placeholders until Phase 7 produces 25-bar, 72-bar, and 200-bar runs.

\section{Benchmark Validation Methodology}
\label{sec:results:validation}
% TODO: state the acceptance tolerance per benchmark (2.0\%, 2.5\%,
% 2.5\%, 3.0\%) as listed in docs/ground\_truth.md; describe seeded
% multi-run protocol and the feasibility criterion.

\section{10-Bar Planar Truss}
\label{sec:results:10bar}

The 10-bar planar truss of Sunar and Belegundu~\cite{sunar1991} is the
most widely reported sizing benchmark in the literature and is used
here as the primary validation case for the classical optimisation
layer. Geometry, loads, material properties, and stress and
displacement limits are reproduced without modification from the
original formulation and are tabulated in
\cref{tab:results:10bar-setup}.

\begin{table}[htbp]
  \centering
  \caption{10-bar planar truss — problem setup (after
    \cite{sunar1991}).}
  \label{tab:results:10bar-setup}
  \begin{tabular}{lll}
    \toprule
    Parameter               & Value                       & Unit \\
    \midrule
    Nodes                   & 6                           & --- \\
    Elements                & 10                          & --- \\
    Young's modulus $E$     & $1.0 \times 10^{7}$         & psi \\
    Density $\rho$          & $0.1$                       & lb/in\textsuperscript{3} \\
    Vertical load at nodes 2, 4 & $100{,}000$             & lbf \\
    Stress limit            & $\pm 25{,}000$              & psi \\
    Displacement limit      & $2.0$                       & in \\
    Area bounds             & $[0.1,\, 35.0]$             & in\textsuperscript{2} \\
    Reference optimum       & $5060.85$                   & lb \\
    \bottomrule
  \end{tabular}
\end{table}

Three algorithms were run against this benchmark with the default
hyperparameters described in \cref{sec:method:optimisers}.
\Cref{tab:results:10bar} reports the best, mean, and standard
deviation of converged weight across independent random seeds, along
with the absolute error of the best solution relative to the published
optimum of $5060.85$~lb. All reported runs are feasible with respect
to both the original stress and displacement limits and the
IS~800:2007 checks described in \cref{sec:method:is800}.

\begin{table}[htbp]
  \centering
  \caption{10-bar planar truss results — best, mean, and standard
    deviation of converged weight across seeds, with percent error
    relative to the literature optimum of \num{5060.85}~lb.}
  \label{tab:results:10bar}
  \sisetup{table-format=4.2}
  \begin{tabular}{lrSSSS}
    \toprule
    Algorithm & Seeds & {Best (lb)} & {Mean (lb)} & {Std.\ (lb)} & {\% err (best)} \\
    \midrule
    GA      & 10 & 5062.78 & 5075.92 & 7.11  & 0.038 \\
    PSO     & 5  & 5061.05 & 5064.70 & 7.21  & 0.004 \\
    NSGA-II & 3  & 5081.51 & 5098.38 & 14.62 & 0.408 \\
    \bottomrule
  \end{tabular}
\end{table}

Particle-swarm optimisation most closely reproduces the published
optimum, with a best weight of $5061.05$~lb corresponding to a
relative error of $0.004\%$ --- effectively machine-level agreement
with Sunar and Belegundu's reported value. The genetic algorithm
converges to within $0.04\%$ across ten independent seeds with a
sample standard deviation of $7.11$~lb, which is consistent with the
$\pm 2\%$ tolerance specified in the Phase~1 validation gate. NSGA-II,
used here in single-objective-weight mode as a smoke-test of the
multi-objective wrapper, returns a slightly heavier best design at
$5081.51$~lb (percent error $0.41\%$) and is primarily exercised later
in \cref{sec:results:pareto} for its multi-objective capabilities.

% TODO (Phase 8): convergence-history figure, area distribution bar
% chart, stress contour plot.
\begin{figure}[htbp]
  \centering
  % \includegraphics[width=0.8\textwidth]{fig_4_2_convergence}
  \caption{Convergence history of GA, PSO, and NSGA-II on the 10-bar
    benchmark. \emph{Figure to be generated in Phase~8.}}
  \label{fig:results:10bar-convergence}
\end{figure}

\section{25-Bar Spatial Truss}
\label{sec:results:25bar}
% TODO (Phase 7 produces runs; Phase 8 writes prose).
% Reference optimum 545.22 lb \cite{schmit1976}, tolerance 2.5\%.

\section{72-Bar Spatial Tower}
\label{sec:results:72bar}
% TODO (Phase 7 runs; Phase 8 prose).
% Reference optimum 379.62 lb \cite{erbatur2000}, tolerance 2.5\%.

\section{200-Bar Planar Truss (Steel)}
\label{sec:results:200bar}
% TODO (Phase 7 runs; Phase 8 prose). Note the steel material-trap:
% $E = 210$ GPa, $\rho = 7850$ kg/m\textsuperscript{3}
% \cite{kaveh2010}.

\section{Neural Surrogate Accuracy and Speedup}
\label{sec:results:surrogate}
% TODO: parity plots, $R^2 = 0.9993$ for weight on 10-bar;
% honest caveat on stress/displacement $R^2 \approx 0.82$/$0.89$;
% wall-time comparison against pure-FEM GA.

\section{PPO Agent Performance}
\label{sec:results:rl}
% TODO: training curves; inference-time advantage; frame as
% inference-speed contribution per Zhao 2021 \cite{zhao2021rl}.

\section{LLM Warm-Start Study}
\label{sec:results:llm}
% TODO: paired t-test, $p = 0.0046$ on 10-bar; 36.8\% reduction in
% generations-to-converge, cache statistics.

\section{Multi-Objective Pareto Fronts}
\label{sec:results:pareto}
% TODO: NSGA-II fronts per benchmark, weight vs max-displacement.

\section{IS 800:2007 Compliance Verification}
\label{sec:results:is800}
% TODO: per-clause pass/fail table for each benchmark's best design;
% hand-calculation spot-check of one critical member per benchmark.
